\chapter{Les composés carbonylés~: additions nucléophiles, acétals et protection}\label{ch:b1:acetals-protection}

Dissolvons du glucose dans l'eau~: presque aucune molécule ne reste sous
la forme d'aldéhyde à chaîne ouverte que dessinent les manuels. La
molécule se replie, son propre groupe hydroxyle s'additionne sur sa
fonction aldéhyde, et un cycle à six atomes se ferme. La même addition
d'un alcool sur un groupe carbonyle, répétée une fois de plus avec un
\omterm{def:b1:elementary-steps:catalyst}{catalyseur} acide, donne les \omterm{def:b1:acetals-protection:hemiacetal}{acétals}, composés stables vis-à-vis des bases
et des \omterm{def:b1:electronic-effects:nucleophile}{nucléophiles} les plus réactifs. Cette stabilité est un outil~: un
aldéhyde qui serait détruit par un \omterm{def:b1:organomagnesium:organometallic}{organomagnésien} peut être caché sous
forme d'\omterm{def:b1:acetals-protection:hemiacetal}{acétal}, traverser la réaction, puis être régénéré à la fin. Ce
chapitre étudie le groupe carbonyle en tant qu'\omterm{def:b1:electronic-effects:nucleophile}{électrophile}, les
additions de l'eau et des alcools, et la stratégie de \omterm{def:b1:acetals-protection:protecting-group}{protection}.

\begin{recall}
Les additions des \omterm{def:b1:organomagnesium:organometallic}{organomagnésiens} sur les composés carbonylés
(\cref{ch:b1:organomagnesium})~; le \omterm{def:b1:extent-q-and-k:quotient}{quotient de réaction} $Q$ et
l'évolution d'un système tant que $Q \neq K$
(\cref{ch:b1:extent-q-and-k})~; \omterm{def:b1:electronic-effects:curly-arrow}{flèches courbes}, \omterm{def:b1:electronic-effects:nucleophile}{nucléophiles} et
\omterm{def:b1:electronic-effects:nucleophile}{électrophiles} (\cref{ch:b1:electronic-effects}). Le volume scolaire
(année 12) a introduit l'idée de protéger un groupe au cours d'une
synthèse.
\end{recall}

\section{Le groupe carbonyle, un électrophile}

La liaison C=O est polarisée vers l'oxygène (\cref{ch:b1:periodicity}),
et une \omterm{def:b1:lewis-resonance-vsepr:resonance}{formule mésomère} \ce{C+-O-} place une charge positive entière sur
le carbone~: celui-ci est \omterm{def:b1:electronic-effects:nucleophile}{électrophile}, attaqué par les \omterm{def:b1:electronic-effects:nucleophile}{nucléophiles}
perpendiculairement au plan du groupe. Les \omterm{def:b1:electronic-effects:nucleophile}{nucléophiles} forts
(\omterm{def:b1:organomagnesium:organometallic}{organomagnésiens}, hydrure, cyanure) s'additionnent directement. Les
\omterm{def:b1:electronic-effects:nucleophile}{nucléophiles} faibles --- eau, alcools --- ont besoin d'aide.

\begin{definition}[Activation électrophile]\label{def:b1:acetals-protection:activation}
L'\emph{activation électrophile}\index{activation electrophile@activation électrophile} d'un
groupe carbonyle est sa protonation (ou sa coordination à un \omterm{def:b1:lewis-resonance-vsepr:lewis-acid}{acide de
Lewis}) sur l'oxygène. Le cation \ce{C=OH+}, dont la charge positive est
partagée par le carbone grâce à la \omterm{def:b1:lewis-resonance-vsepr:resonance}{formule mésomère} \ce{C+-OH}, est
attaqué même par des \omterm{def:b1:electronic-effects:nucleophile}{nucléophiles} faibles.
\end{definition}

\begin{omfigure}
\setchemfig{atom sep=2.1em}
\schemestart
\chemfig{C(-[:150]R)(-[:210]H)=O}
\arrow{->[\ce{H+}]}
\chemfig{C(-[:150]R)(-[:210]H)=O^{+}-H}
\arrow{<->}
\chemfig{C^{+}(-[:150]R)(-[:210]H)-O-H}
\schemestop

{\small Activation d'un aldéhyde en milieu acide~: la protonation de
l'oxygène produit un cation dont une \omterm{def:b1:lewis-resonance-vsepr:resonance}{formule mésomère} porte la charge sur
le carbone.}
\end{omfigure}

\section{Hydrates et hémiacétals}

L'eau s'additionne de façon renversable sur les aldéhydes et les
cétones, \ce{R2C=O + H2O <=>
R2C(OH)2}~; l'hydrate est en général minoritaire pour les cétones et peut
être majoritaire pour le méthanal. Un alcool s'additionne de la même
façon.

\begin{definition}[Hémiacétal et acétal]\label{def:b1:acetals-protection:hemiacetal}
Un \emph{hémiacétal}\index{hemiacetal@hémiacétal} est un composé dont un
carbone porte à la fois un groupe OH et un groupe OR~; il se forme par
addition d'un alcool sur un groupe carbonyle. Un
\emph{acétal}\index{acetal@acétal} est un composé dont un carbone porte
deux groupes OR~; il se forme à partir d'un hémiacétal et d'une seconde
molécule d'alcool, avec perte d'eau.
\end{definition}

Les \omterm{def:b1:acetals-protection:hemiacetal}{hémiacétals} à chaîne ouverte sont en général minoritaires dans leur
équilibre avec l'aldéhyde et l'alcool~; quand le OH et le C=O
appartiennent à la même molécule et qu'un cycle à cinq ou six atomes
peut se fermer, l'\omterm{def:b1:acetals-protection:hemiacetal}{hémiacétal} cyclique prédomine. C'est le cas du glucose,
et du 5-hydroxypentanal, plus simple.

\begin{omfigure}
\setchemfig{atom sep=2em}
\schemestart
\chemfig{HO-[:30]-[:-30]-[:30]-[:-30]-[:30]C(=[:90]O)-[:-30]H}
\arrow{<=>}
\chemfig{*6(O-(-[:-30]OH)-----)}
\schemestop

\medskip

{\small Le 5-hydroxypentanal et son \omterm{def:b1:acetals-protection:hemiacetal}{hémiacétal} cyclique, un cycle à six
atomes contenant l'oxygène. Le OH du carbone 5 s'est additionné sur
l'aldéhyde du carbone 1~; la forme cyclique prédomine, comme pour le
glucose.}
\end{omfigure}

\section{Les acétals}

\begin{proposition}[Mécanisme de formation des acétals]\label{prop:b1:acetals-protection:acetal-mechanism}
En milieu acide, un aldéhyde ou une cétone et un alcool donnent
l'\omterm{def:b1:acetals-protection:hemiacetal}{hémiacétal}, puis l'\omterm{def:b1:acetals-protection:hemiacetal}{acétal}, selon les étapes suivantes~: protonation de
C=O~; addition de l'alcool~; perte d'un proton (\omterm{def:b1:acetals-protection:hemiacetal}{hémiacétal})~;
protonation du OH~; perte d'eau, qui donne un \omterm{def:b1:electronic-effects:intermediates}{carbocation} stabilisé par
le groupe OR restant (un ion oxocarbénium)~; addition d'une seconde
molécule d'alcool~; perte d'un proton. Toutes les étapes sont
renversables.
\end{proposition}

\begin{proof}
Chaque étape est un transfert de proton acido-basique ou l'addition (ou
la perte) d'un \omterm{def:b1:electronic-effects:nucleophile}{nucléophile} sur (ou depuis) un carbone \omterm{def:b1:electronic-effects:nucleophile}{électrophile}, comme
décrit au \cref{ch:b1:electronic-effects}. La perte d'eau est possible
car le cation formé, $\mathrm{R_2C^+{-}OR}$, possède une \omterm{def:b1:lewis-resonance-vsepr:resonance}{formule mésomère} $\mathrm{R_2C{=}O^+R}$
dans laquelle chaque atome a un octet (\cref{prop:b1:electronic-effects:carbocation-order}).
L'acide est consommé lors des protonations et régénéré lors des
déprotonations~: c'est un \omterm{def:b1:elementary-steps:catalyst}{catalyseur}.
\end{proof}

\begin{omfigure}
\setchemfig{atom sep=1.9em}
\schemestart
\chemfig{C(-[:150]R)(-[:210]H)=O}
\arrow{->[\ce{H+}][]}
\chemfig{C^{+}(-[:150]R)(-[:210]H)-OH}
\arrow{->[$\mathrm{R'OH}$][\ce{-H+}]}
\chemfig{C(-[:150]R)(-[:210]H)(-[2]OR')-OH}
\schemestop

\medskip
\schemestart
\chemfig{C(-[:150]R)(-[:210]H)(-[2]OR')-OH}
\arrow{->[\ce{H+}][\ce{-H2O}]}
\chemfig{C^{+}(-[:150]R)(-[:210]H)-OR'}
\arrow{->[$\mathrm{R'OH}$][\ce{-H+}]}
\chemfig{C(-[:150]R)(-[:210]H)(-[2]OR')-OR'}
\schemestop
\par\vspace{6pt}

{\small La formation d'un \omterm{def:b1:acetals-protection:hemiacetal}{acétal}, en deux lignes~: d'abord l'\omterm{def:b1:acetals-protection:hemiacetal}{hémiacétal}
(activation, addition de l'alcool, perte d'un proton), puis l'\omterm{def:b1:acetals-protection:hemiacetal}{acétal}
(protonation du OH, perte d'eau conduisant au cation stabilisé, addition
d'une seconde molécule d'alcool, perte d'un proton).}
\end{omfigure}

La réaction globale, $\mathrm{RCHO} + 2\,\mathrm{R'OH} \rightleftharpoons \mathrm{RCH(OR')_2} + \ce{H2O}$, a
une constante d'équilibre modeste, et l'eau en est un produit.

\begin{proposition}[Déplacer l'acétalisation]\label{prop:b1:acetals-protection:dean-stark}
Si l'eau formée est éliminée en continu, l'acétalisation se poursuit
jusqu'à épuisement du composé carbonylé (ou de l'alcool).
\end{proposition}

\begin{proof}
Le \omterm{def:b1:extent-q-and-k:quotient}{quotient de réaction} $Q = [\text{acétal}][\ce{H2O}]/([\text{aldéhyde}]
[\mathrm{R'OH}]^2)$ reste inférieur à $K$ tant que l'eau est éliminée~:
d'après le \cref{ch:b1:extent-q-and-k}, la réaction continue d'évoluer
dans le sens direct, quelle que soit la valeur de $K$. Avec un diol
(l'éthane-1,2-diol), l'\omterm{def:b1:acetals-protection:hemiacetal}{acétal} est cyclique, et l'équilibre est en outre
plus favorable, une molécule de diol remplaçant deux molécules d'alcool.
\end{proof}

\begin{omfigure}
\begin{tikzpicture}[font=\small]
  \pic at (0,0) {heatingmantle};
  \pic at (0,0) {rbflask={0.4}{omDef!14}};
  \pic (d) at (0,3.0) {deanstark={0.35}{omDef!35}{orange!15}};
  \pic at (0,6.4) {condenser};
  \draw[->] (2.4,4.0) node[right, align=left] {collecteur~: l'eau (phase\\inférieure) ne revient pas~; le\\toluène déborde vers le ballon} -- (0.8,3.8);
  \draw[->] (-2.2,1.0) node[left, align=right] {aldéhyde, diol, toluène,\\catalyseur acide, ébullition} -- (-0.7,0.9);
  \draw[->] (1.8,8.0) node[right] {réfrigérant} -- (0.4,7.8);
\end{tikzpicture}

{\small Un appareil de Dean--Stark. Les vapeurs de toluène et d'eau se
condensent et tombent dans le collecteur gradué~; l'eau, plus dense et
non \omterm{def:b1:intermolecular-forces-solvents:miscibility}{miscible} au toluène, coule au fond et y reste, tandis que le toluène
déborde et retourne dans le ballon. Le volume d'eau recueilli mesure
l'\omterm{def:b1:extent-q-and-k:extent}{avancement} de la réaction.}
\end{omfigure}

\begin{proposition}[Stabilité des acétals]\label{prop:b1:acetals-protection:acetal-stability}
Les \omterm{def:b1:acetals-protection:hemiacetal}{acétals} sont stables en milieu neutre et basique, ainsi que
vis-à-vis des \omterm{def:b1:electronic-effects:nucleophile}{nucléophiles}, des \omterm{def:b1:organomagnesium:organometallic}{organomagnésiens} et des donneurs
d'hydrure~; une solution aqueuse acide les hydrolyse en composé
carbonylé et alcool.
\end{proposition}

\begin{proof}
Le carbone d'un \omterm{def:b1:acetals-protection:hemiacetal}{acétal} ne porte aucun \omterm{def:b1:electronic-effects:nucleophile}{groupe partant} qu'une base ou un
\omterm{def:b1:electronic-effects:nucleophile}{nucléophile} pourrait chasser~: \ce{RO-}, \omterm{def:b1:predominant-reaction:strong-acid}{base forte}, ne part pas, et il
ne reste aucun C=O à attaquer. En solution aqueuse acide, le mécanisme
ci-dessus se déroule en sens inverse~: le large excès d'eau rend $Q < K$
pour la réaction inverse.
\end{proof}

\section{Protéger et déprotéger}

\begin{definition}[Groupe protecteur]\label{def:b1:acetals-protection:protecting-group}
Un \emph{groupe protecteur}\index{groupe protecteur} est un groupe
introduit dans une molécule pour masquer temporairement une fonction,
afin qu'elle ne réagisse pas lors d'une étape visant une autre partie de
la molécule. Son introduction est la \emph{protection}\index{protection},
son retrait la \emph{déprotection}\index{deprotection@déprotection}.
\end{definition}

\begin{method}[Protéger, faire réagir, déprotéger]\label{met:b1:acetals-protection:protect}
\begin{enumerate}
  \item Repérer la fonction qui réagirait, ou qui détruirait le réactif,
        lors de l'étape prévue.
  \item Choisir un \omterm{def:b1:acetals-protection:protecting-group}{groupe protecteur} stable lors de cette étape et
        amovible dans des conditions que le reste de la molécule tolère
        (pour un aldéhyde ou une cétone face à des bases, des
        \omterm{def:b1:electronic-effects:nucleophile}{nucléophiles} ou des hydrures~: un \omterm{def:b1:acetals-protection:hemiacetal}{acétal} cyclique).
  \item Protéger~; réaliser l'étape prévue~; déprotéger.
  \item Faire le compte~: deux étapes de plus, chacune avec son
        \omterm{def:b1:extent-q-and-k:fractional-extent}{rendement}.
\end{enumerate}
\end{method}

\begin{example}[Une cétone qui gêne]\label{ex:b1:acetals-protection:ketoester}
Pour additionner un \omterm{def:b1:organomagnesium:organometallic}{organomagnésien} sur la fonction aldéhyde d'une
molécule qui porte aussi une cétone, aucun choix simple d'\omterm{def:b1:acetals-protection:hemiacetal}{acétal} n'existe
(les aldéhydes forment des \omterm{def:b1:acetals-protection:hemiacetal}{acétals} plus facilement que les cétones, si
bien que ce serait le mauvais groupe qui serait protégé)~; il faut
changer l'ordre des étapes ou les réactifs. La \omterm{def:b1:acetals-protection:protecting-group}{protection} est une
stratégie, pas un réflexe.
\end{example}

\section{Exercices}

\begin{exercise}[$\star$]\label{exo:b1:acetals-protection:1}
Repérer les carbones d'\omterm{def:b1:acetals-protection:hemiacetal}{hémiacétal}, d'\omterm{def:b1:acetals-protection:hemiacetal}{acétal}, d'hydrate, d'éther ou
d'alcool dans~:
\ce{CH3CH(OH)OCH3}, \ce{CH3CH(OCH3)2}, \ce{CH3CH(OH)2}, \ce{CH3OCH3},
\ce{(CH3)2C(OCH2CH3)2}.
\end{exercise}

\begin{exercise}[$\star$]\label{exo:b1:acetals-protection:2}
Écrire l'équation globale de la formation de l'\omterm{def:b1:acetals-protection:hemiacetal}{acétal} du propanal avec
le méthanol, et de l'\omterm{def:b1:acetals-protection:hemiacetal}{acétal} cyclique de la propanone avec
l'éthane-1,2-diol.
\end{exercise}

\begin{exercise}[$\star$]\label{exo:b1:acetals-protection:3}
Lesquels de ces réactifs laissent un \omterm{def:b1:acetals-protection:hemiacetal}{acétal} inchangé~: une solution
d'hydroxyde de sodium, le bromure de méthylmagnésium, l'acide sulfurique
dilué dans l'eau, le borohydrure de sodium~?
\end{exercise}

\begin{exercise}[$\star$]\label{exo:b1:acetals-protection:4}
Dessiner l'\omterm{def:b1:acetals-protection:hemiacetal}{hémiacétal} cyclique formé par le 4-hydroxybutanal. Quelle est
la taille du cycle~?
\end{exercise}

\begin{exercise}[$\star\star$]\label{exo:b1:acetals-protection:5}
Écrire le mécanisme complet, avec des \omterm{def:b1:electronic-effects:curly-arrow}{flèches courbes}, de la formation
acido-catalysée de l'\omterm{def:b1:acetals-protection:hemiacetal}{acétal} diméthylique de l'éthanal, et montrer que
l'acide est régénéré.
\end{exercise}

\begin{exercise}[$\star\star$]\label{exo:b1:acetals-protection:6}
Écrire le mécanisme de l'hydrolyse acide du
2,2-diméthoxypropane. Pourquoi utilise-t-on une grande quantité d'eau~?
\end{exercise}

\begin{exercise}[$\star\star$]\label{exo:b1:acetals-protection:7}
On suit l'acétalisation de \qty{0.250}{mol} d'une cétone avec un
appareil de Dean--Stark. Quel volume d'eau doit-on recueillir à
conversion totale (masse volumique voisine de \qty{1.0}{g/mL})~? Au bout
d'une heure, on a recueilli
\qty{3.2}{mL}~: quel est le taux de conversion~?
\end{exercise}

\begin{exercise}[$\star\star$]\label{exo:b1:acetals-protection:8}
Expliquer pourquoi un \omterm{def:b1:acetals-protection:hemiacetal}{acétal} cyclique formé avec l'éthane-1,2-diol se
forme plus complètement que l'\omterm{def:b1:acetals-protection:hemiacetal}{acétal} formé avec deux molécules de
méthanol, pour un même composé carbonylé.
\end{exercise}

\begin{exercise}[$\star\star$]\label{exo:b1:acetals-protection:9}
Pourquoi faut-il éliminer (ou neutraliser) le \omterm{def:b1:elementary-steps:catalyst}{catalyseur} acide avant de
traiter un composé protégé par un \omterm{def:b1:organomagnesium:organometallic}{organomagnésien}~?
\end{exercise}

\begin{exercise}[$\star\star\star$]\label{exo:b1:acetals-protection:10}
Proposer une synthèse de la 5-hydroxy-5-méthylhexan-2-one,
\ce{(CH3)2C(OH)CH2CH2COCH3}, à partir de la 4-chlorobutan-2-one et de la
propanone, en utilisant une \omterm{def:b1:acetals-protection:protecting-group}{protection}. Donner chaque étape, ses
réactifs et le rôle de chacun.
\end{exercise}

\begin{exercise}[$\star\star\star$]\label{exo:b1:acetals-protection:11}
Montrer à l'aide du \omterm{def:b1:extent-q-and-k:quotient}{quotient de réaction} qu'un large excès de méthanol
déplace lui aussi l'acétalisation d'un aldéhyde, et comparer cette
méthode à l'élimination de l'eau.
\end{exercise}

\begin{exercise}[$\star\star\star$]\label{exo:b1:acetals-protection:12}
Expliquer pourquoi l'\omterm{def:b1:acetals-protection:hemiacetal}{hémiacétal} cyclique du glucose réagit, dans le
méthanol acide, pour donner un \omterm{def:b1:acetals-protection:hemiacetal}{acétal} méthylique (un méthylglycoside)
alors que les autres groupes OH du glucose ne sont pas convertis en
éthers.
\end{exercise}

\section{Problème~: un organomagnésien qui porte un aldéhyde}

\begin{problem}[{Problème du week-end --- pourquoi le 4-bromobutanal ne
peut pas donner d'organomagnésien, sa protection sous forme d'acétal
cyclique avec un appareil de Dean--Stark, l'addition de l'organomagnésien
sur la propanone, la déprotection, et le volume d'eau recueilli}]
\label{pb:b1:acetals-protection:1}
Une chimiste a besoin de 5-hydroxy-5-méthylhexanal,
\ce{(CH3)2C(OH)CH2CH2CH2CHO}, et prévoit de le préparer à partir du
4-bromobutanal, \ce{BrCH2CH2CH2CHO}, et de la propanone.
Elle part de \qty{15.09}{g} de 4-bromobutanal. Masses molaires
(\unit{g/mol})~: \ce{C4H7BrO} 150.9, \ce{C2H6O2} 62.0, \ce{C6H11BrO2} 194.9,
\ce{C3H6O} 58.0, \ce{C9H18O3} 174.0, \ce{C7H14O2} 130.0, \ce{H2O} 18.0~;
masse volumique de l'eau \qty{0.997}{g/mL}.
% ledger: rho:water

\medskip
\noindent\textbf{Partie I --- Le problème.}
\begin{enumerate}
  \item Écrire l'\omterm{def:b1:organomagnesium:organometallic}{organomagnésien} que donnerait le 4-bromobutanal.
  \item Pourquoi ce réactif ne peut-il pas exister~? Écrire la réaction
        qui le détruit.
  \item Quelle fonction faut-il protéger, et contre quoi~?
  \item Pourquoi un \omterm{def:b1:acetals-protection:hemiacetal}{acétal} convient-il~?
  \item Pourquoi préfère-t-on un \omterm{def:b1:acetals-protection:hemiacetal}{acétal} cyclique (avec
        l'éthane-1,2-diol)~?
  \item Calculer la quantité de matière de 4-bromobutanal.
  \item Calculer la masse de diol pour un excès de \qty{20}{\percent}.
\end{enumerate}

\medskip
\noindent\textbf{Partie II --- La \omterm{def:b1:acetals-protection:protecting-group}{protection}.}
\begin{enumerate}[resume]
  \item Écrire l'équation de l'acétalisation avec l'éthane-1,2-diol.
  \item Quel est le rôle du \omterm{def:b1:elementary-steps:catalyst}{catalyseur} acide (l'acide
        4-méthylbenzènesulfonique)~?
  \item Décrire le fonctionnement de l'appareil de Dean--Stark.
  \item Expliquer à l'aide de $Q$ et $K$ pourquoi l'élimination de l'eau
        déplace la réaction.
  \item Pourquoi est-ce le toluène, et non l'eau, qui retourne dans le
        ballon~?
  \item Comment la chimiste sait-elle que la réaction est terminée~?
  \item Calculer la masse de bromure protégé attendue à conversion
        totale.
\end{enumerate}

\medskip
\noindent\textbf{Partie III --- L'étape de Grignard.}
\begin{enumerate}[resume]
  \item Écrire la formation de l'\omterm{def:b1:organomagnesium:organometallic}{organomagnésien} à partir du bromure
        protégé.
  \item Écrire son addition sur la propanone et l'\omterm{def:b1:alcohol-activation:alkoxide}{alcoolate} formé.
  \item Pourquoi l'hydrolyse de cette étape doit-elle être faiblement
        basique ou neutre (par exemple avec une solution aqueuse de
        chlorure d'ammonium) plutôt qu'acide~?
  \item Calculer la masse d'alcool protégé pour un \omterm{def:b1:extent-q-and-k:fractional-extent}{rendement} de
        \qty{70}{\percent} sur cette étape.
\end{enumerate}

\medskip
\noindent\textbf{Partie IV --- La \omterm{def:b1:acetals-protection:protecting-group}{déprotection}.}
\begin{enumerate}[resume]
  \item Écrire la \omterm{def:b1:acetals-protection:protecting-group}{déprotection} en solution aqueuse acide.
  \item Pourquoi l'alcool tertiaire résiste-t-il à une solution aqueuse
        faiblement acide~?
  \item Le produit peut fermer un cycle à six atomes par addition de son
        OH sur l'aldéhyde. Dessiner cet \omterm{def:b1:acetals-protection:hemiacetal}{hémiacétal} cyclique.
  \item Calculer la masse d'aldéhyde-alcool obtenue pour un \omterm{def:b1:extent-q-and-k:fractional-extent}{rendement} de
        \omterm{def:b1:acetals-protection:protecting-group}{déprotection} de \qty{90}{\percent}.
  \item Quel \omterm{def:b1:extent-q-and-k:fractional-extent}{rendement} global la chimiste atteint-elle à partir du
        4-bromobutanal~?
  \item Calculer la masse d'eau formée lors de l'étape de \omterm{def:b1:acetals-protection:protecting-group}{protection} à
        conversion totale.
  \item Donner le volume d'eau recueilli dans l'appareil de Dean--Stark à
        conversion totale.
\end{enumerate}
\end{problem}
